\label{subsec:analytic_geometry_parametric_line_equations}

\begin{DefinitionBox}{Parameterdarstellung einer Geraden}{analytic_geometry_parametric_line}
  Eine Gerade $g$ mit Aufpunktvektor $\overrightarrow{OP}$ und Richtungsvektor
  $\vec a$ besitzt die
  \emph{Parameterdarstellung} \EnglishTerm{parametric equation}.

  \noindent
  \begin{minipage}[t]{0.54\linewidth}
    \vspace{0pt}
    \centering
    \CachedTikzFigure{content/resources/chapters/analytic_geometry/tikz/parametric_line_equation_annotation.tex}
  \end{minipage}\hfill
  \begin{minipage}[t]{0.42\linewidth}
    \vspace{0pt}
    \centering
    \begingroup
      \def\LineUnit{0.75}
      \def\LineSupportX{3}
      \def\LineSupportY{3}
      \def\LineDirectionX{1}
      \def\LineDirectionY{-1}
      \def\LineLambdaMin{-3.75}
      \def\LineLambdaMax{3.75}
      \CachedTikzFigure{content/resources/chapters/analytic_geometry/tikz/parameterized_line_representation.tex}
    \endgroup
  \end{minipage}

  $\lambda\in\R$ wirkt dabei wie ein Schieberegler, d.h. wenn
  $\lambda$ alle reellen Werte durchläuft, zeigt der Ortsvektor
  $\vec x=\overrightarrow{OX}$ auf alle Punkte der Geraden $g$.
\end{DefinitionBox}

\begin{ExampleBox}{Punkte einer Geraden bestimmen}{analytic_geometry_line_points_example}
\begin{minipage}[t]{0.50\linewidth}
  \vspace{0pt}
  Wir betrachten die folgende Gerade in Parameterdarstellung:
  \[
    g:\quad
    \vec x=\begin{pmatrix}2\\2\end{pmatrix}
    +\lambda\begin{pmatrix}1\\-1\end{pmatrix}
  \]
  Für drei verschiedene Parameterwerte bestimmen wir die zugehörigen Ortsvektoren:
  \[
    \begin{aligned}
      \color{red!75!black}
      \lambda=+1:\quad
      \vec x_{+1}
        &=\begin{ClozeContinuation}[10]{2}\begin{pmatrix}2\\2\end{pmatrix}
          +(+1)\begin{pmatrix}1\\-1\end{pmatrix}
        =\begin{pmatrix}3\\1\end{pmatrix}\end{ClozeContinuation},\\[0.9em]
      \color{teal!80!black}
      \lambda=+3:\quad
      \vec x_{+3}
        &=\begin{ClozeContinuation}[10]{2}\begin{pmatrix}2\\2\end{pmatrix}
          +3\begin{pmatrix}1\\-1\end{pmatrix}
        =\begin{pmatrix}5\\-1\end{pmatrix}\end{ClozeContinuation},\\[0.9em]
      \color{violet!80!black}
      \lambda=-1:\quad
      \vec x_{-1}
        &=\begin{ClozeContinuation}[10]{2}\begin{pmatrix}2\\2\end{pmatrix}
          +(-1)\begin{pmatrix}1\\-1\end{pmatrix}
        =\begin{pmatrix}1\\3\end{pmatrix}\end{ClozeContinuation}.
    \end{aligned}
  \]
Alle drei Ortsvektoren zeigen auf Punkte derselben Geraden.
\end{minipage}\hfill
\begin{minipage}[t]{0.46\linewidth}
  \vspace{0pt}
  \centering
  \CachedTikzFigure{content/resources/chapters/analytic_geometry/tikz/line_parameter_points.tex}
\end{minipage}
\par\smallskip
\end{ExampleBox}
